% =====================================================================
%  Préambule — Medicacom AOI Transformation
%  Compilation : pdflatex (deux passes) ou latexmk -pdf main.tex
% =====================================================================

% ---------------------------------------------------------- encodage & langue
\usepackage[T1]{fontenc}
\usepackage[utf8]{inputenc}
\usepackage[french]{babel}
\usepackage{lmodern}
\usepackage{amsmath}
\usepackage{amssymb}
\usepackage{textcomp}
\usepackage{microtype}

% ------------------------------------------------------------------ mise en page
\usepackage[a4paper,left=2.6cm,right=2.4cm,top=2.8cm,bottom=2.6cm]{geometry}
\usepackage{setspace}
\onehalfspacing

% --------------------------------------------------------------------- couleurs
\usepackage[table]{xcolor}
\definecolor{medicaDark}{HTML}{10231C}   % vert sombre de la charte
\definecolor{medicaGreen}{HTML}{157F5C}  % vert accent
\definecolor{medicaDeep}{HTML}{0F5E44}   % vert foncé
\definecolor{medicaTeal}{HTML}{1C6B7D}   % teal (données)
\definecolor{medicaAmber}{HTML}{A66A16}  % ambre (attention)
\definecolor{medicaRed}{HTML}{B23A3A}    % rouge (risque)
\definecolor{medicaSoft}{HTML}{E6F1EB}   % vert très clair
\definecolor{medicaBg}{HTML}{EDF2EF}     % fond
\definecolor{medicaBorder}{HTML}{E1E8E4} % bordure
\definecolor{medicaMuted}{HTML}{5E6F67}  % texte atténué

% ------------------------------------------------------------------ graphiques
\usepackage{graphicx}
\usepackage{float}
\usepackage{tikz}
\usetikzlibrary{
  shapes.geometric, shapes.misc, arrows.meta, positioning, fit, calc,
  backgrounds, decorations.pathreplacing, patterns, matrix, chains
}
\usepackage{pgfplots}
\pgfplotsset{compat=1.18}

% --------------------------------------------------------------------- tableaux
\usepackage{booktabs}
\usepackage{array}
\usepackage{longtable}
\usepackage{tabularx}
\usepackage{multirow}
\usepackage{makecell}
\renewcommand{\arraystretch}{1.28}
\newcolumntype{L}[1]{>{\raggedright\arraybackslash}p{#1}}
\newcolumntype{C}[1]{>{\centering\arraybackslash}p{#1}}
\newcolumntype{R}[1]{>{\raggedleft\arraybackslash}p{#1}}
\newcolumntype{Y}{>{\raggedright\arraybackslash}X}

% ----------------------------------------------------------------------- listes
\usepackage{enumitem}
\setlist{itemsep=2pt, topsep=4pt, parsep=0pt}
\setlist[itemize]{label=\textcolor{medicaGreen}{\textbullet}}

% -------------------------------------------------------------------- légendes
\usepackage{caption}
\usepackage{subcaption}
\captionsetup{
  font=small, labelfont={bf,color=medicaDeep},
  textfont=it, width=0.92\textwidth, skip=7pt
}

% ------------------------------------------------------------------- en-têtes
\usepackage{fancyhdr}
\pagestyle{fancy}
\fancyhf{}
\fancyhead[L]{\small\textcolor{medicaMuted}{\nouppercase{\leftmark}}}
\fancyhead[R]{\small\textcolor{medicaMuted}{Medicacom AOI Transformation}}
\fancyfoot[C]{\small\thepage}
\renewcommand{\headrulewidth}{0.4pt}
\renewcommand{\footrulewidth}{0pt}
\fancypagestyle{plain}{\fancyhf{}\fancyfoot[C]{\small\thepage}%
  \renewcommand{\headrulewidth}{0pt}}

% --------------------------------------------------------------------- titres
\usepackage{titlesec}
\titleformat{\chapter}[display]
  {\normalfont\huge\bfseries\color{medicaDeep}}
  {\filright\Large\textcolor{medicaGreen}{\chaptertitlename\ \thechapter}}
  {8pt}{\Huge}
\titlespacing*{\chapter}{0pt}{-14pt}{26pt}
\titleformat{\section}{\Large\bfseries\color{medicaDeep}}{\thesection}{0.7em}{}
\titleformat{\subsection}{\large\bfseries\color{medicaDark}}{\thesubsection}{0.7em}{}
\titleformat{\subsubsection}{\normalsize\bfseries\color{medicaDark}}{\thesubsubsection}{0.7em}{}

% ------------------------------------------------------------------- encadrés
\usepackage{framed}

% Encadré « Présentation orale suggérée » : sécable sur plusieurs pages.
\newenvironment{oral}{%
  \par\medskip
  \def\FrameCommand{\textcolor{medicaGreen}{\vrule width 3pt}\hspace{12pt}}%
  \MakeFramed{\advance\hsize-\width\FrameRestore}%
  \small\textbf{\textcolor{medicaDeep}{Présentation orale suggérée}}\par\smallskip
  \itshape\ignorespaces
}{\endMakeFramed\par\medskip}

% Encadré neutre pour une remarque importante.
\newenvironment{remarque}[1][Remarque]{%
  \par\medskip
  \def\FrameCommand{\textcolor{medicaTeal}{\vrule width 3pt}\hspace{12pt}}%
  \MakeFramed{\advance\hsize-\width\FrameRestore}%
  \small\textbf{\textcolor{medicaTeal}{#1}}\par\smallskip\ignorespaces
}{\endMakeFramed\par\medskip}

% Encadré d'avertissement, pour les écarts entre l'attendu et le réalisé.
\newenvironment{ecart}[1][Écart avec le cahier des charges]{%
  \par\medskip
  \def\FrameCommand{\textcolor{medicaAmber}{\vrule width 3pt}\hspace{12pt}}%
  \MakeFramed{\advance\hsize-\width\FrameRestore}%
  \small\textbf{\textcolor{medicaAmber}{#1}}\par\smallskip\ignorespaces
}{\endMakeFramed\par\medskip}

% ------------------------------------------------------------------ code source
\usepackage{listings}
\lstdefinestyle{medicacode}{
  basicstyle=\ttfamily\scriptsize,
  keywordstyle=\color{medicaDeep}\bfseries,
  commentstyle=\color{medicaMuted}\itshape,
  stringstyle=\color{medicaTeal},
  numberstyle=\tiny\color{medicaMuted},
  numbers=left, numbersep=8pt,
  backgroundcolor=\color{medicaBg},
  frame=leftline, rulecolor=\color{medicaBorder}, framesep=8pt,
  breaklines=true, breakatwhitespace=false,
  showstringspaces=false, tabsize=2,
  captionpos=b, xleftmargin=14pt,
  extendedchars=true, literate=%
    {é}{{\'e}}1 {è}{{\`e}}1 {ê}{{\^e}}1 {ë}{{\"e}}1
    {à}{{\`a}}1 {â}{{\^a}}1 {î}{{\^\i}}1 {ï}{{\"\i}}1
    {ô}{{\^o}}1 {ö}{{\"o}}1 {û}{{\^u}}1 {ù}{{\`u}}1 {ü}{{\"u}}1
    {ç}{{\c c}}1 {É}{{\'E}}1 {È}{{\`E}}1 {À}{{\`A}}1 {Ç}{{\c C}}1
    {€}{{\texteuro}}1 {×}{{$\times$}}1 {→}{{$\rightarrow$}}1 {−}{{$-$}}1
    {«}{{\guillemotleft}}1 {»}{{\guillemotright}}1 {’}{{'}}1 {…}{{...}}3
}
\lstset{style=medicacode}
\lstdefinelanguage{TypeScript}{
  keywords={const,let,var,function,return,if,else,for,while,import,export,
            from,type,interface,extends,implements,class,new,async,await,
            default,as,of,in,null,undefined,true,false},
  sensitive=true,
  comment=[l]{//}, morecomment=[s]{/*}{*/},
  morestring=[b]", morestring=[b]'
}

% ----------------------------------------------------------------- hyperliens
\usepackage[hidelinks,
  pdftitle={Conception et développement de Medicacom AOI Transformation},
  pdfsubject={Plateforme SaaS de pilotage de l'adoption de l'IA dans l'industrie pharmaceutique},
  pdfkeywords={IA, TOGAF, readiness, portefeuille, pharmaceutique, FastAPI, React}
]{hyperref}
\usepackage[french]{cleveref}

% =====================================================================
%  Commandes utilitaires
% =====================================================================

% Emplacement de capture d'écran. Affiche le vrai fichier s'il existe dans
% images/, sinon un cadre explicite — le document compile dans les deux cas.
\newlength{\captureL}
\newlength{\captureH}
\newcommand{\capture}[3]{%
  \setlength{\captureL}{#3\textwidth}%
  \setlength{\captureH}{0.56\captureL}%
  \begin{figure}[H]
  \centering
  \IfFileExists{images/#1}{%
    \includegraphics[width=\captureL]{images/#1}%
  }{%
    \begin{tikzpicture}
      \draw[dashed, color=medicaBorder, line width=1pt, fill=medicaBg]
        (0pt,0pt) rectangle (\captureL, \captureH);
      \node[align=center, text=medicaMuted, font=\small]
        at (0.5\captureL, 0.5\captureH)
        {\textbf{Capture à insérer}\\[3pt]\ttfamily images/#1};
    \end{tikzpicture}%
  }
  \caption{#2}
  \label{fig:#1}
  \end{figure}%
}

% Marqueurs d'état d'implémentation, employés dans les tableaux de synthèse.
\newcommand{\fait}{\textcolor{medicaGreen}{\textbf{Implémenté}}}
\newcommand{\partiel}{\textcolor{medicaAmber}{\textbf{Partiel}}}
\newcommand{\absent}{\textcolor{medicaRed}{\textbf{Non implémenté}}}
\newcommand{\oui}{\textcolor{medicaGreen}{\textbf{\checkmark}}}
\newcommand{\non}{\textcolor{medicaMuted}{--}}

% Raccourcis de nommage
\newcommand{\appname}{Medicacom AOI Transformation}
\newcommand{\code}[1]{\texttt{\small #1}}

% Styles TikZ partagés par les diagrammes
\tikzset{
  bloc/.style={
    rectangle, rounded corners=4pt, draw=medicaBorder, fill=white,
    text width=2.9cm, align=center, minimum height=1cm, font=\small
  },
  blocvert/.style={bloc, fill=medicaSoft, draw=medicaGreen},
  blocsombre/.style={bloc, fill=medicaDark, draw=medicaDark, text=white},
  blocteal/.style={bloc, fill=white, draw=medicaTeal},
  fleche/.style={-{Stealth[length=5pt]}, draw=medicaMuted, thick},
  flechevert/.style={-{Stealth[length=5pt]}, draw=medicaGreen, thick},
  acteur/.style={
    rectangle, draw=medicaDeep, fill=medicaSoft, rounded corners=3pt,
    text width=2.5cm, align=center, font=\scriptsize\bfseries, minimum height=0.8cm
  },
  cas/.style={
    ellipse, draw=medicaBorder, fill=white, align=center,
    font=\scriptsize, minimum width=2.9cm, minimum height=0.82cm,
    inner sep=2pt
  },
  classe/.style={
    rectangle, draw=medicaDeep, fill=white, align=left,
    font=\scriptsize, inner sep=0pt
  },
  lien/.style={draw=medicaMuted, thin},
}

% En-tête de classe UML : nom + liste d'attributs.
\newcommand{\umlclasse}[2]{%
  \begin{tabular}{@{}l@{}}
    \cellcolor{medicaSoft}\textbf{\strut~#1~}\\
    \hline
    \begin{minipage}{\dimexpr\linewidth}\vspace{2pt}#2\vspace{2pt}\end{minipage}
  \end{tabular}%
}
